\section{Modul 00. Register, interupsi, dan kerangka program}

\subsection{Tujuan}

Setelah menyelesaikan modul ini, mahasiswa mampu:
\begin{enumerate}[nosep]
    \item membaca bilangan heksadesimal dan kode ASCII yang dipakai
          pada layanan DOS;
    \item menjelaskan pembagian register \texttt{AX} dan \texttt{DX}
          beserta tugas setiap bagiannya;
    \item menjelaskan cara kerja interupsi perangkat lunak dan peran
          \texttt{INT 21h};
    \item membedakan kerangka program \texttt{.COM} dan \texttt{.EXE};
    \item membaca program contoh dan menelusuri isi register tanpa
          menulis atau merakit program.
\end{enumerate}

\subsection{Dasar singkat}

\subsubsection*{Bilangan heksadesimal dan kode ASCII}

Bilangan heksadesimal memakai basis 16 dengan digit 0 sampai 9 dan
A sampai F. Pada assembler TASM, bilangan heksadesimal ditulis
dengan akhiran \texttt{h}. Bila digit pertama berupa huruf, angka nol
harus ditambahkan di depan agar assembler tidak membacanya sebagai
label; karena itu ditulis \texttt{0Dh}, bukan \texttt{Dh}. Nilai yang
dipakai pada modul ini dirangkum pada Tabel~\ref{tab:m00-hex}.

\begin{table}[h]
    \centering
    \caption{Nilai heksadesimal yang dipakai pada Modul 00.}
    \label{tab:m00-hex}
    \small
    \begin{tabular}{@{}lll@{}}
        \toprule
        \textbf{Heksadesimal} & \textbf{Desimal} & \textbf{Arti} \\
        \midrule
        \texttt{02h} & 2 & nomor layanan cetak satu karakter \\
        \texttt{09h} & 9 & nomor layanan cetak string \\
        \texttt{0Ah} & 10 & \emph{line feed}, turun satu baris \\
        \texttt{0Dh} & 13 & \emph{carriage return}, kembali ke awal baris \\
        \texttt{24h} & 36 & tanda \texttt{\$} penutup string \\
        \bottomrule
    \end{tabular}
\end{table}

Karakter disimpan sebagai kode ASCII satu byte. Huruf \texttt{A}
bernilai \texttt{41h}, angka \texttt{0} bernilai \texttt{30h}, dan
spasi bernilai \texttt{20h}. Kode \texttt{0Dh} mengembalikan kursor
ke awal baris, sedangkan \texttt{0Ah} menurunkannya satu baris.
Karena itu pasangan \texttt{0Dh}, \texttt{0Ah} dipakai untuk
berpindah ke awal baris berikutnya. Tanda \texttt{\$} bukan bagian
kalimat yang dicetak; layanan \texttt{AH=09h} memakainya sebagai
penanda akhir string.

\subsubsection*{Register AX, DX, dan register lain}

Mikroprosesor 8086 bekerja dengan register 16 bit. Register
\texttt{AX} dapat dipisah menjadi dua bagian 8 bit, yaitu
\texttt{AH} (byte tinggi) dan \texttt{AL} (byte rendah). Register
\texttt{DX} juga terbagi menjadi \texttt{DH} dan \texttt{DL}.
Sebagai contoh, bila \texttt{AX} bernilai \texttt{1234h}, maka
\texttt{AH} bernilai \texttt{12h} dan \texttt{AL} bernilai
\texttt{34h}. Peran bagian register yang dipakai layanan DOS
ditunjukkan pada Tabel~\ref{tab:m00-register}.

\begin{table}[h]
    \centering
    \caption{Bagian register yang dipakai layanan DOS.}
    \label{tab:m00-register}
    \small
    \begin{tabular}{@{}ll@{}}
        \toprule
        \textbf{Register} & \textbf{Isi atau tugas pada modul ini} \\
        \midrule
        \texttt{AH} & nomor layanan yang dipilih sebelum \texttt{INT 21h} \\
        \texttt{AL} & karakter masukan (dibahas pada Modul 02) \\
        \texttt{DL} & satu karakter yang dicetak layanan \texttt{02h} \\
        \texttt{DX} & offset string yang dicetak layanan \texttt{09h} \\
        \texttt{DS} & nomor segmen tempat data program berada \\
        \bottomrule
    \end{tabular}
\end{table}

Register lain cukup dikenal namanya. \texttt{BX} menyimpan alamat
offset; \texttt{CX} menjadi pencacah pengulangan; \texttt{SI}
menunjuk alamat sumber; \texttt{IP} menyimpan offset instruksi
berikutnya; \texttt{CS} menunjuk segmen kode; \texttt{SS} menunjuk
segmen stack; \texttt{ES} menunjuk segmen tambahan; dan
\texttt{FLAGS} mencatat keadaan hasil operasi terakhir.

\subsubsection*{Interupsi perangkat lunak}

Interupsi perangkat lunak adalah permintaan dari program kepada
prosesor untuk mengerjakan suatu layanan. Ketika interupsi terjadi,
program berhenti sejenak, prosesor menjalankan rutin layanan, lalu
program dilanjutkan pada instruksi berikutnya. Prosesor 8086
menyediakan 256 nomor interupsi, yaitu 0 sampai 255. Nomor dipilih
dengan instruksi \texttt{INT}.

\texttt{INT 21h} adalah interupsi layanan DOS. Nomor fungsi
diletakkan pada \texttt{AH}, sedangkan argumen lain diletakkan pada
register sesuai fungsi. Fungsi yang diperkenalkan pada modul ini
adalah \texttt{02h}, \texttt{09h}, dan \texttt{4Ch}. Khusus program
\texttt{.COM}, program dapat ditutup dengan \texttt{INT 20h}.
Rangkuman fungsi tersebut ditunjukkan pada Tabel~\ref{tab:m00-int21}.

\begin{table}[h]
    \centering
    \caption{Fungsi \texttt{INT 21h} dan \texttt{INT 20h} pada Modul 00.}
    \label{tab:m00-int21}
    \small
    \begin{tabular}{@{}lll@{}}
        \toprule
        \textbf{Perintah} & \textbf{Fungsi} & \textbf{Register} \\
        \midrule
        \texttt{INT 21h}, \texttt{AH=02h} & cetak satu karakter
        & \texttt{DL} \\
        \texttt{INT 21h}, \texttt{AH=09h} & cetak string berakhir
        \texttt{\$} & \texttt{DS:DX} \\
        \texttt{INT 21h}, \texttt{AH=4Ch} & akhiri program
        & \texttt{AL} = kode keluar \\
        \texttt{INT 20h} & akhiri program \texttt{.COM} & -- \\
        \bottomrule
    \end{tabular}
\end{table}

\subsubsection*{Kerangka program .COM dan .EXE}

Program \texttt{.COM} memakai model \texttt{TINY} dan direktif
\texttt{ORG 100h}. Seluruh kode dan data berada dalam satu segmen,
sehingga \texttt{DS} tidak perlu diisi sendiri. Data diletakkan di
dalam segmen kode dan dilewati dengan pola \texttt{JMP MULAI} agar
tidak ikut dieksekusi. Program ditutup dengan \texttt{INT 20h}.

\begin{lstlisting}[numbers=none]
.MODEL TINY
.CODE
ORG 100h
START:
    JMP MULAI           ; lompati data
    ; data diletakkan di sini
MULAI:
    ; instruksi program
    INT 20h             ; akhiri program .COM
END START
\end{lstlisting}

Program \texttt{.EXE} memakai model \texttt{SMALL}, menyediakan
stack, dan memisahkan segmen data dari segmen kode. Sebelum data
disentuh, \texttt{DS} harus diisi dengan alamat segmen data melalui
\texttt{AX}. Program ditutup dengan \texttt{AH=4Ch} dan
\texttt{INT 21h}.

\begin{lstlisting}[numbers=none]
.MODEL SMALL
.STACK 100h
.DATA
    ; variabel diletakkan di sini
.CODE
START:
    MOV AX, @DATA       ; alamat segmen data
    MOV DS, AX          ; pasang ke DS
    ; instruksi program
    MOV AH, 4Ch         ; akhiri program .EXE
    INT 21h
END START
\end{lstlisting}

Direktif \texttt{ORG 100h} hanya dipakai pada program \texttt{.COM}
dan tidak boleh dicampur dengan kerangka \texttt{.EXE}.

Rujukan singkat: Gunawan, Bagian 3 dan 4; S'to, Bab I, III, IV, V,
dan VI; buku ajar, \texttt{chapters/03}, \texttt{05}, \texttt{07},
dan \texttt{09/section-03.tex}.

\subsection{Alat dan bahan}

\begin{itemize}[nosep]
    \item Komputer dengan emu8086 untuk menampilkan isi register.
    \item DOSBox dengan DEBUG sebagai alat pemeriksa.
    \item TASM dan TLINK sebagai perakit dan pengait; diperkenalkan
          pada modul ini, tetapi belum dijalankan.
    \item Berkas \texttt{code/m00-contoh-int21.asm} dalam bentuk
          cetak atau digital.
    \item Lembar jurnal praktikum.
\end{itemize}

\subsection{Contoh program}

Program berikut adalah program \texttt{.COM} lengkap. Program ini
mencetak satu kalimat dengan layanan \texttt{AH=09h}, mencetak satu
karakter dengan layanan \texttt{AH=02h}, lalu berhenti dengan
\texttt{INT 20h}. Perhatikan pola \texttt{JMP MULAI}: data
\texttt{PESAN} diletakkan di antara instruksi lompatan dan label
\texttt{MULAI} agar tidak dieksekusi sebagai instruksi.

\lstinputlisting[
    caption={Contoh program \texttt{.COM} dengan \texttt{INT 21h}.},
    label={lst:m00-contoh}
]{code/m00-contoh-int21.asm}

Pada program ini \texttt{DS} tidak diisi karena \texttt{.COM}
memakai satu segmen. Sebelum \texttt{INT 21h} pertama,
\texttt{AH} berisi \texttt{09h} dan \texttt{DX} berisi offset
\texttt{PESAN}. Sebelum \texttt{INT 21h} kedua, \texttt{AH} berisi
\texttt{02h} dan \texttt{DL} berisi karakter \texttt{'!'}.

\subsection{Langkah kerja}

\begin{enumerate}
    \item Baca program contoh. Tandai tiga bagian: direktif,
          data, dan instruksi. Tuliskan nomor baris setiap bagian.
    \item Sebelum \texttt{INT 21h} pertama, catat isi \texttt{AH} dan
          \texttt{DX}. Jelaskan ke mana \texttt{DX} menunjuk.
    \item Jelaskan peran \texttt{0Dh} dan \texttt{0Ah} di dalam
          \texttt{PESAN}, lalu jelaskan apa yang terjadi bila tanda
          \texttt{\$} dihapus.
    \item Sebelum \texttt{INT 21h} kedua, catat isi \texttt{AH} dan
          \texttt{DL}. Tentukan karakter yang muncul di layar.
    \item Bandingkan kerangka \texttt{.COM} dan \texttt{.EXE} pada
          dasar singkat. Isi tabel: model memori, penggunaan
          \texttt{ORG 100h}, inisialisasi \texttt{DS}, dan perintah
          penutup.
    \item Jelaskan mengapa \texttt{JMP MULAI} diperlukan dan apa
          akibatnya bila data diletakkan tanpa lompatan.
\end{enumerate}

\subsection{Tugas}

\begin{enumerate}
    \item Diberikan potongan berikut, tentukan fungsi yang dipanggil,
          register yang dipakai, dan hasilnya:
          \texttt{MOV AH,09h}; \texttt{MOV DX,OFFSET PESAN};
          \texttt{INT 21h}.
    \item Jelaskan mengapa program \texttt{.COM} tidak memerlukan
          \texttt{MOV AX,@DATA} dan \texttt{MOV DS,AX}.
    \item Sebutkan tiga perbedaan kerangka program \texttt{.COM} dan
          \texttt{.EXE}.
    \item Terjemahkan nilai \texttt{02h}, \texttt{09h},
          \texttt{0Dh}, \texttt{0Ah}, dan \texttt{24h} ke desimal,
          lalu tuliskan artinya.
    \item Jelaskan apa yang terjadi bila \texttt{INT 20h} dihilangkan
          dari program \texttt{.COM}.
\end{enumerate}

\subsection{Hasil yang dikumpulkan}

\begin{enumerate}
    \item Lembar jurnal berisi jawaban langkah kerja dan tugas.
    \item Tabel isi register \texttt{AH}, \texttt{AL}, \texttt{DX},
          \texttt{DL}, dan \texttt{DS} sebelum serta sesudah setiap
          \texttt{INT}.
    \item Tabel perbandingan kerangka \texttt{.COM} dan \texttt{.EXE}.
    \item Catatan singkat tentang peran \texttt{0Dh}, \texttt{0Ah},
          dan tanda \texttt{\$}.
\end{enumerate}
